\subsection{From One Bar to the Rest of the Session}
\label{sec:multihorizon}
\label{sec:multihorizon_results}

The forecasts of Table~\ref{tab:block_ladder} are one-bar forecasts.
This subsection extends them to the variance remaining in the session,
Equation~\eqref{eq:mh_target}, using the constructions of
Section~\ref{sec:multihorizon_methods}.

% PARKED 2026-08-19: horizon-profile results + tab:multih_profile held out
\iffalse
\paragraph{The horizon profile.} Table~\ref{tab:multih_profile} refits
the incumbent and the two-block ridge directly at five horizons from
thirty minutes to a trading week, on one shared evaluation span of
\mhpNBars{} bars per horizon. Both models lose accuracy as the horizon
lengthens, as they must. What changes across the rows is the
\emph{difference} between them. At one bar the two arms are level on
this span (paired DM \mhpHOneBlkDM{}); at two hours the ridge is behind
(\mhpHFourBlkDM{}); at eight hours it is level again
(\mhpHSixteenBlkDM{}); at one day it is ahead at
\mhpHFortyEightBlkDM{}; and at one week it is ahead by
\mhpHTwoFortyBlkRel{}\% of the incumbent's loss, at
\mhpHTwoFortyBlkDM{}. The exogenous block's advantage is not a property
of the thirty-minute horizon that decays away from it: it is smallest
there and largest at the longest horizon fit. The information in the
wide block is about slow-moving variance, and it pays most where the
HAR lag structure has the least to say. The two-hour row is the one
exception in sign and is not significant; it is reported as found.

\begin{table}[H]
\centering
\caption{Direct forecasts at longer horizons. Both arms are refit at
each horizon $h$ with the paper's estimator and penalties, on the same
final \mhpNBars{} bars of the panel; the target is the realized
variance of bar $t+h-1$. QLIKE under the evaluation protocol
(Section~\ref{sec:smearing_contract}); $\Delta$ is the two-block ridge
minus the incumbent on common rows; DM~$t$ is the paired
Diebold--Mariano statistic, negative favouring the ridge. The one-bar
row is a subsample of Table~\ref{tab:block_ladder}'s panel and is
reported here only to anchor the profile.
\label{tab:multih_profile}}
\begin{tabular}{@{}lcccrr@{}}
\toprule
Horizon & $n$ & OLS--HAR & Two-block & $\Delta$QLIKE & DM~$t$ \\
\midrule
1 bar (30 min) & 57{,}200 & 0.2422 & 0.2422 & $-$0.0000 & $-$0.0 \\
4 bars (2 h) & 57{,}200 & 0.2778 & 0.2811 & +0.0033 & +1.7 \\
16 bars (8 h) & 57{,}200 & 0.3040 & 0.3025 & $-$0.0014 & $-$0.8 \\
48 bars (1 day) & 57{,}200 & 0.3289 & 0.3200 & $-$0.0088 & $-$3.0 \\
240 bars (5 days) & 57{,}200 & 0.3951 & 0.3365 & $-$0.0586 & $-$14.9 \\
\bottomrule

\end{tabular}
\end{table}

\fi

\paragraph{The rest of the session.} The object is the variance
remaining before the close, Equation~\eqref{eq:mh_target}, forecast at
every bar from the one-step forecast standing then, on \mhNDays{}
complete regular sessions. Two questions are asked of it: how much of
the remaining session a single next-bar forecast reaches, and what the
running realization of the current session adds; and whether the
two-block ridge's one-bar advantage survives the extension to a horizon
of hours.

\begin{table}[H]
\centering
\caption{Multi-horizon forecasts of the variance remaining in the
session, pooled over entry hours 10:00--15:00~ET (one loss per session,
the mean over the eleven entry hours). QLIKE on remaining variance,
lower is better. DM~$t$ (ridge vs.\ incumbent) is the paired
Diebold--Mariano statistic of the two-block ridge against the OLS--HAR
incumbent under the same construction; DM~$t$ (vs.\ benchmark) is the
two-block ridge against the horizon-specific unconditional benchmark.
Negative favours the ridge. All constructions use only information
available at the entry bar.
\label{tab:multihorizon_pooled}}
\begin{tabular}{@{}llcccc@{}}
\toprule
Construction & Sample & OLS--HAR & Two-block & DM~$t$ (ridge vs.\ inc.) & DM~$t$ (vs.\ bench.) \\
\midrule
Horizon benchmark (no model) & All sessions & 0.8409 & --- & --- & --- \\
One-step forecast, mapped to the horizon (M1) & All sessions & 0.1837 & 0.1740 & $-$4.2 & $-$19.9 \\
\quad $+$ realized so far this session (M2) & All sessions & 0.1814 & 0.1719 & $-$5.4 & $-$19.8 \\
\quad $+$ previous session, same horizon (M3) & All sessions & 0.1761 & 0.1674 & $-$5.9 & $-$20.0 \\
M2, one regression with hour dummies & All sessions & 0.1814 & 0.1719 & $-$5.3 & $-$19.8 \\
\midrule
Horizon benchmark (no model) & 2016--2024 & 0.9044 & --- & --- & --- \\
One-step forecast, mapped to the horizon (M1) & 2016--2024 & 0.2537 & 0.2357 & $-$3.2 & $-$13.6 \\
\quad $+$ realized so far this session (M2) & 2016--2024 & 0.2511 & 0.2346 & $-$3.9 & $-$13.4 \\
\quad $+$ previous session, same horizon (M3) & 2016--2024 & 0.2404 & 0.2264 & $-$4.0 & $-$13.7 \\
M2, one regression with hour dummies & 2016--2024 & 0.2506 & 0.2341 & $-$3.9 & $-$13.5 \\
\bottomrule

\end{tabular}
\end{table}

\paragraph{How far one bar reaches.} Mapped alone to the rest of the
session (M1), the incumbent's next-bar forecast scores QLIKE
\mhAzeroMOneQLIKE{} against the horizon benchmark's \mhNaiveQLIKE{},
and the two-block ridge's scores \mhBlkMOneQLIKE{}: a single next-bar
forecast carries substantial information about the hours that follow
(DM \mhBlkMOneDMVsNaive{} against the benchmark for the ridge). Adding
the variance already realized in the session (M2) improves both models
further --- the ridge to \mhBlkMTwoQLIKE{} --- and yesterday's
realization of the same horizon (M3) improves them again, to
\mhBlkMThreeQLIKE{} for the ridge. Pooling the eleven entry hours into
one regression with hour dummies scores the same as the eleven
hour-specific fits (DM \mhPooledVsPerHourBlk{} for the ridge), even
though the slope from ``the next bar'' to ``the rest of the session''
moves through the day (\mhBlkSlopeOnestepEarly{} at 10:00 against
\mhBlkSlopeOnestepLate{} at 14:30 on the one-step term, with the
running realization taking up the difference: \mhBlkSlopeSofarEarly{}
against \mhBlkSlopeSofarLate{}; descriptive, full sample).

\paragraph{The ridge's advantage at the longer horizon.} Under every
construction the two-block ridge scores below the incumbent. Mapped
alone, the paired DM is \mhMOneDMBlkVsAzero{}; with the running
realization, \mhMTwoDMBlkVsAzero{}; with the seasonal anchor,
\mhMThreeDMBlkVsAzero{}. Table~\ref{tab:multihorizon_byhour} gives the
M2 comparison by entry hour: the ridge scores below the incumbent at
\mhHoursBlkBetter{} of the \mhHoursTotal{} entry hours, with paired DM
ranging from \mhDMBlkVsAzeroMin{} to \mhDMBlkVsAzeroMax{}. The one-bar
advantage of Table~\ref{tab:block_ladder} therefore carries to the
remaining-session horizon: the exogenous block improves the forecast
an intraday user needs at every bar, not only the next thirty minutes.
The 2016--2024 rows of Table~\ref{tab:multihorizon_pooled} show the
same ordering.

\begin{table}[H]
\centering
\caption{Multi-horizon forecasts by entry hour, construction M2 (next-bar
forecast plus realized variance so far), all sessions. QLIKE on the
variance remaining from that hour to the close; DM~$t$ as in
Table~\ref{tab:multihorizon_pooled}. The remaining horizon shrinks from
eleven bars at 10:00 to one at 15:00.
\label{tab:multihorizon_byhour}}
\begin{tabular}{@{}lccccc@{}}
\toprule
Entry & Benchmark & OLS--HAR & Two-block & DM~$t$ (ridge vs.\ inc.) & DM~$t$ (vs.\ bench.) \\
\midrule
10:00 & 0.7024 & 0.1330 & 0.1257 & $-$4.9 & $-$19.1 \\
10:30 & 0.7302 & 0.1408 & 0.1338 & $-$4.0 & $-$19.4 \\
11:00 & 0.7589 & 0.1507 & 0.1436 & $-$4.2 & $-$19.4 \\
11:30 & 0.7860 & 0.1632 & 0.1548 & $-$3.7 & $-$19.3 \\
12:00 & 0.8121 & 0.1830 & 0.1732 & $-$4.4 & $-$18.9 \\
12:30 & 0.8382 & 0.2005 & 0.1879 & $-$3.9 & $-$18.7 \\
13:00 & 0.8615 & 0.2182 & 0.2095 & $-$2.8 & $-$18.4 \\
13:30 & 0.8914 & 0.2410 & 0.2276 & $-$3.6 & $-$18.3 \\
14:00 & 0.9248 & 0.2835 & 0.2731 & $-$3.8 & $-$17.5 \\
14:30 & 0.9557 & 0.1492 & 0.1378 & $-$6.1 & $-$21.6 \\
15:00 & 0.9889 & 0.1327 & 0.1238 & $-$6.3 & $-$19.2 \\
\bottomrule

\end{tabular}
\end{table}


\paragraph{The per-bar path, refit directly.}
Table~\ref{tab:dense_byk} refits both arms at each of the eleven bars
beyond the one Table~\ref{tab:block_ladder} scores, on \dmhN{} origins
across the panel. The horizon convention matters for reading it: the
paper's one-bar forecast pairs the information standing at the start of
a bar with that bar's variance; the first row here targets the bar
after that, so the table covers the second through the twelfth bar
ahead --- the horizons between Table~\ref{tab:block_ladder} and a full
session, at bar-level resolution. On every one of them the exogenous
block does not help: the two-block ridge scores above the incumbent
from DM \dmhKOneDM{} at the second bar to \dmhKElevenDM{} at the
twelfth (regular-hours origins \dmhRthKOneDM{} to
\dmhRthKElevenDM{}), in every era of the panel separately, and under
per-bar refit as well as the daily refit reported in the table (a
per-bar-refit diagnostic on the first eighth of the panel narrows the
gap materially --- the cadence is not free for a wide, heavily shrunk
block --- but does not change its sign). The one-bar advantage of
Table~\ref{tab:block_ladder} therefore does not extend to the hours
immediately beyond it: the exogenous information is about the next
thirty minutes and, as the rest-of-session results above show in
aggregate, about the level of the session --- not about the particular
bar two hours from now.

\begin{table}[H]
\centering
\caption{Direct per-bar forecasts at bar-level horizons beyond
Table~\ref{tab:block_ladder}: the row labelled $t+k$ targets the
$(k+1)$-th bar ahead of the information set, so the table covers the
second through the twelfth bar. Both arms are refit at every horizon
with the paper's design and penalties, coefficients re-solved daily.
QLIKE under the evaluation protocol on all origins from the first full
window; $\Delta$ is the two-block ridge minus the incumbent on common
rows; DM~$t$ is the paired Diebold--Mariano statistic, negative
favouring the ridge.
\label{tab:dense_byk}}
\begin{tabular}{@{}lcccrr@{}}
\toprule
Target bar & $n$ & OLS--HAR & Two-block & $\Delta$QLIKE & DM~$t$ \\
\midrule
$t+1$ & 273{,}371 & 0.1928 & 0.1979 & +0.0051 & +11.2 \\
$t+2$ & 273{,}371 & 0.1988 & 0.2061 & +0.0073 & +14.6 \\
$t+3$ & 273{,}371 & 0.2027 & 0.2109 & +0.0082 & +15.9 \\
$t+4$ & 273{,}371 & 0.2057 & 0.2152 & +0.0096 & +18.9 \\
$t+5$ & 273{,}371 & 0.2081 & 0.2183 & +0.0102 & +19.4 \\
$t+6$ & 273{,}371 & 0.2099 & 0.2210 & +0.0111 & +20.5 \\
$t+7$ & 273{,}371 & 0.2112 & 0.2229 & +0.0117 & +20.9 \\
$t+8$ & 273{,}371 & 0.2124 & 0.2245 & +0.0121 & +21.0 \\
$t+9$ & 273{,}371 & 0.2136 & 0.2258 & +0.0122 & +20.8 \\
$t+10$ & 273{,}371 & 0.2152 & 0.2274 & +0.0122 & +20.4 \\
$t+11$ & 273{,}371 & 0.2160 & 0.2283 & +0.0123 & +20.2 \\
\bottomrule

\end{tabular}
\end{table}
